\section{在边界上办事的人}

\subsection{馆伴的行囊}

辽宋使节的相遇，通常由盟书、国书和皇帝的称谓开始；真正让相遇发生的，却是一支必须带着衣物、食物、印信、翻译和随从穿越关津的队伍。馆伴要知道哪一处渡口可用，哪位地方官负责供应，哪一种礼仪会在抵达时被要求。他的差事不是单纯传话：一项称谓写错、一次马匹供应延误、一次随从越过规定区域，都可能被中枢解释为对盟约的试探。澶渊以后，\eventref{LIAO-1005-EVENT-038}{盟约的落实}正是在这些看起来琐碎的事务里被反复确认。

宋人使辽行录保留了许多馆舍、饮食、服饰和仪注，因而读来像一扇可以打开辽社会的窗。窗的框架却由使者自己的任务决定。他尤其注意座次、仪仗和让自己感到陌生的语言，却未必记录替他添草料、修车轮和传递文书的人。辽方材料更少留下同一行程的镜像。把两类沉默并置后，读者可以把使团当作一项基础设施来理解：它需要国家授权，也需要沿途城镇愿意承担日常成本；它可显示政权关系的温度，不能替代普通边民的完整声音。

西夏使节面对的道路更不稳定。兴庆府、陕西沿边、河西与辽境之间的通道，会因战事、封市、洪水或地方冲突而改变。使者携带的礼物和文书不仅是王权的象征，也是让接待者能够把这支队伍识别为“可交涉的人”的凭证。宋方记录中的请封、朝贡或谢恩，常把西夏使节置于本朝礼制；西夏的年号、官署与文字则表明它并不只在宋朝的词典里行动。我们很难从现存材料复原一趟行程的每个夜宿点，但能看见外交必须依赖一批熟悉地名、语言和危险的人。

\subsection{关市的早晨}

开市前，关市并不是一片自然形成的热闹。守门者要确认日子和禁令，商人要判断带去的是布帛、茶、盐、皮毛还是马匹，地方官要处理税、检查和纠纷。辽宋边市的和平价值不在于取消边界，而在于把相遇压进可以预期的时间和地点。马匹、铁器和消息的流动仍会引起军事担忧；一处市场同时是收入来源、情报节点和潜在风险。对携货的人而言，一道临时禁令可能比远方朝廷的盟约更直接地改变一年生计。

西夏与宋的互市也不应被写成贡赐制度的附属物。河套、陇右和河西的居民在城寨、牧地和绿洲之间交换物品，官方允许、禁止或征税的对象，会影响哪些交换能留下文书。史书往往在战争、禁市和和议时才让市场进入叙事，平常日子里的小额交易则几乎无声。钱币、窖藏和遗址能补出一些物的移动，却不能替任何一位商人说明价格、利润或风险。市场史的困难正是它的价值：它迫使我们把国号与实际生活之间的关津、度量和信用写出来。

市场也带来语言问题。讨价还价、核验货物、解释禁令和解决债务，都需要有人能在不同语词之间使意思暂时对得上。翻译者可能是官员、商人、僧侣、地方首领或同时承担几种身份的人。材料很少让他们以自己的口吻出现，然而没有他们，条约中的“互市”无法变成可执行的安排。将多语世界只写成几种文字并列，会错过这些口头工作；将它写成毫无摩擦的交流，又会忘记误译、欺诈和权力不均也会沿同一条道路移动。

\subsection{北面与南面之间的账簿}

辽朝的北面、南面官制常被用来说明一个草原政权如何统治农耕地区。这个解释若停在名称上，反而会把办事的人抹去。官署需要能够登记、传达、裁判和催纳的书吏；军队需要知道粮草何时到达；地方豪强、寺院和部众首领会用自己掌握的资源影响命令的落点。北面与南面不是两只内部完全一致的盒子，而是不同地区、不同来源的资源被朝廷尝试接合的方式。

一份制度条文可以告诉我们官名和职掌，却不能证明每一处官署都按同样方法工作。辽的百官志、营卫志和食货志本身也是后出的整理，它们把长期变化压缩为较稳定的制度图像。契丹文字题记、汉文墓志和城址材料能将某些官名、家族和地点重新钉在地面上；它们又偏向能刻石、能安葬、能被发掘的少数人。读者不应在“中央文献”与“地方材料”之间选择一方为真，而应让二者互相限制：前者说明制度想要什么，后者偶尔说明谁能够把它变成事实。

西夏的官文书和法律材料提出了相似的问题。法卷可以规定如何分类财物、罪名和程序，实际的案件却需要当事人、书写者和裁断者都进入同一套手续。黑水城保存的材料使这种书写世界可见，却不能让我们推断全境都有同样密集的官署与识字者。西夏本国连续国史未存，王朝政治的年代常要由宋、辽、金、元史互证；正因如此，一份出土文书的局部性不是缺陷，而是防止敌国叙事独占解释权的必要边界。

\subsection{一座寺院如何穿过战争}

寺院的屋顶、经卷和造像看上去远离边防，实际上也要面对道路和资源。修一座塔需要砖木与工匠，抄一部经需要纸、墨和能读写的人，维持僧众则需要食物、土地或捐施。辽代的寺院题记、经幢和墓葬图像显示佛教如何同家族记忆、城市声望和政治保护相连；它们不证明全体居民共享同一种信仰，也不说明国家能无条件支配寺院。宗教网络之所以值得进入编年，是因为它提供了战争、迁徙和政权更替之外仍在运作的资源与文字通道。

西夏佛经印本尤其提醒人们，宗教不是只有口头仪式的领域。翻译、雕版、抄写和供养要把多种语言、技术与资金聚在一起。汉文、藏文、西夏文和回鹘文文本在同一地区出现，不等于每个居民都能读写每一种语言；它更明确地告诉我们，寺院、官署与商路需要处理不同读者和不同权威。某一部经书的保存，说明有人曾使它成为值得保存的物，不说明它代表西夏所有地区或所有宗派。

战争会让寺院的保护关系改变，却未必让宗教生活立即停止。辽亡后，原有寺院可能在金的官府、地方家族或新的施主网络中寻找继续存在的方法；西夏在蒙古征服后也有文本、图像和人员以不同方式进入新的政治环境。征服史书关心城破和君主投降，往往少写一处寺院何时重修、一批经卷怎样转移。不能由这种沉默虚构连续性，但也不能把王朝终局写成所有书写与祈愿同时消失。

\subsection{墓地里的迁徙}

辽代汉文墓志和契丹文字材料常为后人安排一个人的祖先、官历、婚姻与葬地。它们有意把家族写成有秩序的延续，追赠的官衔和理想化的品行不能直接当作生前事实；然而，正因为要解释一个人为何属于某个家族，迁徙和任职的路线也会露出来。祖籍、任所和埋葬地不在一处，提示道路、军役、婚姻和官署能把家庭拉向不同的中心。

西夏墓葬和题记的保存更不均衡，不能拼出同样连续的家族地图。黑水城文书中偶见的姓名、债务、施主或官署关系，能够把抽象的“党项社会”拆成若干具体当事人，却很少告诉我们其一生如何度过。女性、儿童、雇佣者、牧人和流动者尤其难以在史料中持续追踪。承认这些人难以被听见，不是把他们逐出叙事；相反，国家征役、修渠、战争运输和市场禁令之所以重要，正因为它们会落到这些没有留下墓志的人身上。

\subsection{女真与青唐：第三方改变所有人的选择}

辽宋关系在一一一四年之前已不是两方独自安排的边界。女真诸部的行动、东北的贡纳与贸易、宋廷获得的消息，都让燕云和河北的计算发生变化。宋人把女真来使写成边政消息，辽廷则要处理原有驻防和部众关系的松动；女真首领也在不同邻居之间寻找资源。后来的金朝兴起不应被倒灌成一种天命，它是许多道路、据点、部众和政治承诺重新组合的结果。

西夏面向的第三方同样不止宋和辽。青唐、河西绿洲、吐蕃诸部与回鹘社群构成不同的贸易和军事空间。西夏向甘凉方向的行动会触及城镇、水源和宗教网络，而不只是扩张一条地图边线。宋方史料常从本朝边防看待这种变化，西夏文书则更容易显示某一城镇的行政与宗教实践；两者之间仍有大量未被保存的地方经验。外部世界不是一张附录地图，而是促使辽、西夏统治者反复调整联盟、驻军和使节的压力来源。

\subsection{河流冻结时}

北方的道路不是全年可用的直线。河流封冻可以让某些路段更易通过，也可能让补给和渡运变得危险；春汛、泥泞、风雪和缺水会把军队、使团和商队的时间表打乱。辽军南下、宋军北伐、西夏的边地调动，都必须处理这一层季节性。环境不是决定胜败的独立主角，同一条河、同一段山路会因准备、地方合作和战争局势而产生不同结果；但任何忽略它的叙事都会把人和物写得像在空地图上移动。

灾异记录必须更谨慎地使用。正史在本纪中记录旱、蝗、洪水和地震，常与政治劝戒相连；一条记录可说明朝廷注意到某种危机，不能自动证明整个政权的人口或财政同样受灾。水利遗址、沉积和地方材料能补充环境史，也有各自的时间尺度。对辽与西夏而言，较可靠的问题不是“自然是否导致衰亡”，而是某次水、路和气候的变化怎样让既有的征发、储粮、迁徙或防守更难维持。
